% =============================================================================
%  VII. DISCUSSION — KEY FINDING
% =============================================================================
\begin{frame}{VII. Discussion --- Key Finding}
The performance collapse reported by Wang et al.~\cite{wang2025srpm} is not
only a model evaluation issue. This study shows that it is strongly related
to how the sepsis label is constructed.

\vspace{6pt}
\textbf{The Sepsis-3 label in MIMIC-IV has three important limitations:}
\begin{enumerate}
  \item \textbf{Depends on clinical actions} --- No sepsis cases are
        identified before treatment starts~\cite{kamran2024evaluation}.
  \item \textbf{Sensitive to label definition} --- Different reasonable
        interpretations can produce almost double the number of sepsis
        cases~\cite{cohen2024subtle}.
  \item \textbf{Produces poor clinical utility} --- All tested models show
        negative Utility Scores~\cite{wang2025srpm}.
\end{enumerate}

\vspace{6pt}
\begin{alertblock}{}
  \centering These findings suggest that improving model complexity alone
  cannot solve the problem. Better sepsis label design is required.
\end{alertblock}
\end{frame}


% =============================================================================
%  VII. DISCUSSION — CONTRIBUTIONS
% =============================================================================
\begin{frame}{VII. Discussion --- Contributions}
\footnotesize\renewcommand{\arraystretch}{1.15}
\begin{center}
\begin{tabular}{cp{11cm}}
\toprule
\textbf{\#} & \textbf{Contribution} \\
\midrule
C1 & Compared three Sepsis-3 label definitions using the same prediction
     model on MIMIC-IV~\cite{cohen2024subtle}. \\
C2 & Showed no sepsis cases exist before antibiotics or culture tests,
     highlighting the label depends on clinical
     actions~\cite{kamran2024evaluation}. \\
C3 & Confirmed all models had negative Utility Scores, extending
     Wang et al.~\cite{wang2025srpm}. \\
C4 & Found 10 clinical features changed their effects across label
     definitions, supporting Lauritsen et al.~\cite{lauritsen2021framing}. \\
C5 & Quantified split optimism: random splits inflated AUROC by 0.123
     vs.\ temporal splits, confirming Guo et al.~\cite{guo2022temporal}. \\
C6 & First study to examine how sepsis label definitions affect racial
     subgroup analysis~\cite{wang2022equity}. \\
C7 & Used a pre-registered, TRIPOD+AI-compliant, fully reproducible
     pipeline~\cite{collins2024tripodai}. \\
\bottomrule
\end{tabular}
\end{center}
\end{frame}


% =============================================================================
%  VII. DISCUSSION — LIMITATIONS
% =============================================================================
\begin{frame}{VII. Discussion --- Limitations}
\begin{itemize}
  \item \textbf{Small external event count:} The eICU dataset had only
        \textbf{63 sepsis cases}, making some evaluation metrics less
        reliable.
  \item \textbf{Deep learning comparison:} A deep survival model was planned
        but not implemented.
  \item \textbf{Limited features:} Clinical notes, procedure codes, and
        ventilator settings were not included.
  \item \textbf{Specific to Sepsis-3:} These findings apply to the Sepsis-3
        label and may not apply to other sepsis definitions.
\end{itemize}

\vspace{8pt}
\begin{exampleblock}{}
  \centering These limits are documented rather than hidden.
\end{exampleblock}
\end{frame}


% =============================================================================
%  VIII. CONCLUSION AND FUTURE WORK
% =============================================================================
\begin{frame}{VIII. Conclusion and Future Work}
\begin{enumerate}
  \item \textbf{Better sepsis labels:} Define sepsis using objective patient
        measurements (such as lactate levels and organ function) instead of
        treatment decisions.
  \item \textbf{Prospective study:} Test the new label in real clinical
        settings using future patient data.
  \item \textbf{Better alert threshold:} Choose the alert threshold that
        maximises clinical benefit (Utility Score) instead of only improving
        prediction accuracy.
  \item \textbf{Multi-hospital evaluation:} Test the model in different
        hospitals and identify an acceptable number of false alerts for
        clinical use.
\end{enumerate}
\end{frame}
